\clearpage
\begingroup
\fancyhead[L]{\small THEORIE DER BEWEGUNG DES MONDES.}
\renewcommand{\headrulewidth}{0pt}
\SourceLeaf{626}
\section*{DRITTER ABSCHNITT.}
\begin{center}
\itshape Erste Annäherung.
\end{center}

\par\noindent\hspace*{\parindent}Das erste, was wir hier zu thun haben, ist die Bestimmung von \(m''\).
Aus der Theorie der elliptischen Bewegung ist bekannt, dass, wenn man die
Bewegung der Erde um die Sonne bloss als Product der Anziehungskräfte dieser
beiden Weltkörper betrachtet, und die mittlere Bewegung der Sonne \(=nu\),
[die] mittlere Entfernung der Erde von der Sonne \(=B\) setzt, sein werde:
\[
m'+m''=nnB^3\frac{\mathrm{d}u^2}{\mathrm{d}t^2};
\]
\noindent also, \(\dfrac{m'}{m''}=\mu\) gesetzt, welches Verhältniss hinreichend genau
bekannt ist \(\left(=\dfrac{1}{365412}\right)\),
\[
m''=nnB^3\frac{\mathrm{d}u^2}{\mathrm{d}t^2}\frac{1}{1+\mu}
\mathpunct{\raisebox{.4ex}{.}}
\]
\noindent Ferner ist
\[
R'=\BareRoot{\left(RR-\frac{2R}{p}\cos(v-V)+\frac{1+\theta\theta}{pp}\right)}
\]
\noindent oder, mit Weglassung der höhern Potenzen von \(\dfrac{1}{R}\) und \(\theta\),
\[
\begin{aligned}
\frac{1}{R'^3}&=\frac{1}{R^3}\left\{1+\frac{3\cos(v-V)}{Rp}\right.\\
&\qquad\left.+\frac{15}{4}\frac{1}{RRpp}(1+\cos2(v-V))
-\frac32\frac{1}{RRpp}\right\}\mathpunct{\raisebox{.4ex}{.}}
\end{aligned}
\]
\noindent Hieraus folgt
\[
\begin{aligned}
\pi&=\frac{nn}{1+\mu}\frac{B^3}{R^3}\left\{\frac32\sin2(v-V)\right.\\
&\qquad\left.+\frac{1}{Rp}\left(\frac38\sin(v-V)+\frac{15}{8}\sin3(v-V)\right)\right\}\\[5pt]
\omega&=\frac{nn}{1+\mu}\frac{B^3}{R^3}\left\{\frac32+\frac32\cos2(v-V)\right.\\
&\qquad\left.+\frac{1}{Rp}\left(\frac{33}{8}\cos(v-V)+\frac{15}{8}\cos3(v-V)\right)\right\}\\[5pt]
\psi&=\frac{nn}{1+\mu}\frac{B^3}{R^3}\left\{\frac12+\frac32\cos2(v-V)\right.\\
&\qquad\left.+\frac{1}{Rp}\left(\frac98\cos(v-V)+\frac{15}{8}\cos3(v-V)\right)\right\}
\mathpunct{\raisebox{.4ex}{.}}
\end{aligned}
\]
\SourceLeaf{627}
\noindent Endlich, wenn man die Sonnenexcentricität \(=\varepsilon\) und die mittlere
Anomalie \(=a\) setzt,
\[
\frac{R}{B}=1+\tfrac12\varepsilon\varepsilon+\varepsilon\cos a
-\tfrac12\varepsilon\varepsilon\cos2a,
\]
\noindent also
\[
\frac{B^3}{R^3}=1+\tfrac32\varepsilon\varepsilon-3\varepsilon\cos a
+\tfrac92\varepsilon\varepsilon\cos2a,
\]
\noindent und
\[
V=nu-2\varepsilon\sin a+\tfrac54\varepsilon\varepsilon\sin2a.
\]

\par\noindent\hspace*{\parindent}Ferner bemerken wir, dass wir bei Substitution der elliptischen Bewegung
in den Theilen \(\sin\lambda\) und \(\cos M\) die Zeichen \(\lambda\), \(M\)
nothwendig unbestimmt beibehalten müssen und nicht \(v-\mathit{Const.}\)
dafür substituiren dürfen. Es könnte nemlich wohl sein, dass die Perturbation
der Sonne dem Knoten und Apogäum eine nicht bloss periodische Bewegung mittheilte,
also \(\lambda\) und \(M\) die Form \(\zeta v-\mathit{Const.}\) bekämen,
so dass \(\zeta\) eine von der Einheit wenig verschiedene Grösse würde.
In diesem Falle würde die Voraussetzung der elliptischen Bewegung durchaus falsch,
wenn man \(\lambda\) anstatt \(=\zeta v-\mathit{Const.}\), \(=v-\mathit{Const.}\)
setzen wollte, weil diese Verschiedenheit sich immer mehr aufhäufen würde.
Wir lassen daher die Argumente \(\lambda\) und \(M\) unbestimmt, oder vielmehr
setzen \(\lambda=\zeta v-\mathit{Const.}\), \(M=\eta v-\mathit{Const.}\),
so dass \(\zeta\), \(\eta\) Coefficienten bedeuten, deren Werth ohne die Perturbation
der Sonne \(=1\) sein würde, jetzt aber so lange unbestimmt bleibt,
bis die Rechnung darüber entscheidet.

\par\noindent\hspace*{\parindent}Setzen wir bei dieser ersten Annäherung in der dritten Gleichung die
Theile, die \(\dfrac{nn}{R}\), \(nn\varepsilon\) und \(nne\) enthalten,
bei Seite, so wird dieselbe
\[
\begin{aligned}
0&=\frac{\mathrm{dd}\theta}{\mathrm{d}v^2}+(1+\tfrac32nn)\theta\\
&\qquad+\tfrac32nn\left(\theta\cos2(v-V)
-\frac{\mathrm{d}\theta}{\mathrm{d}v}\sin2(v-V)\right);
\end{aligned}
\]
\noindent statt \(V\) dürfen wir hier \(nv[+\mathit{Const.}]\) setzen,
folglich wird, wenn wir statt \(v-nv[-\mathit{Const.}]\), \(E\) und statt
\(\theta\), \(g\sin\lambda\) schreiben,
\[
0=\frac{\mathrm{dd}\theta}{\mathrm{d}v^2}+(1+\tfrac32nn)\theta
-\tfrac32nng\sin(2E-\lambda).
\]

\par\noindent\hspace*{\parindent}Diese Gleichung zeigt nun sogleich, dass nicht mehr
\(\lambda=v-\mathit{Const.}\), sondern
\(=v\BareRoot{(1+\tfrac32nn)}-\mathit{Const.}\) zu setzen sei;
also \(\dfrac{\mathrm{d}\lambda}{\mathrm{d}v}=1+\tfrac34nn\), oder dass der
Knoten eine rückwärtsgehende Bewegung habe, sie sich zur mittlern Bewegung
des Mondes wie \(\tfrac34nn:1\) verhalte. Hieraus folgt also der erste genäherte
\SourceLeaf{628}Werth
\[
\theta=g\sin\lambda+(\tfrac38n+\tfrac{3}{32}nn)g\sin(2E-\lambda).
\]

\par\noindent\hspace*{\parindent}Um doch zu sehen, was für numerische Coefficienten diese erste
Annäherung gibt, bemerken wir, dass den Beobachtungen zufolge
\[
\begin{array}{r@{\;}c@{\;}r}
\log n&=&8{,}873\,9096\\
\log nn&=&7{,}747\,8192\\
\log\tfrac34&=&9{,}875\,0613\\\cline{2-3}
&&7{,}622\,8805
\end{array}
\]
\noindent also \(\tfrac34nn=0{,}004\,1964\,35\), da die Beobachtungen geben:
\(\dfrac{\mathrm{d}\lambda}{\mathrm{d}v}=1{,}004\,0218\,706\).

\par\noindent\hspace*{\parindent}Ferner ist \(g\) in Secunden \(=18567\overset{\prime\prime}{,}933\):
\[
\begin{array}{r@{\;}c@{\;}r@{\qquad}r@{\;}c@{\;}r}
\log&=&4{,}268\,7636&&&4{,}268\,7636\\
\log\tfrac38n&=&8{,}447\,9409&\log\tfrac{3}{32}nn&=&6{,}719\,7905\\\cline{2-3}\cline{5-6}
&&2{,}716\,7045&&&0{,}988\,5541\\[3pt]
&&520\overset{\prime\prime}{,}840&&&+9\overset{\prime\prime}{,}740
\end{array}
\]
\noindent also der Coefficient \(=530\overset{\prime\prime}{,}580\), welches nach einer
in der Folge zu erklärenden Reduction genau mit den Beobachtungen übereinstimmt.
\clearpage
\endgroup
